% =============================================================================
% Pillar 3 -- Iter 111: log-G linear fit + best-G scalability, G=4 vs G=32
% Driver: scripts/group_size_analysis.py (iter111 path)
% Inputs:  experiments/results/groupsize_zvf_sweep.tsv
%          experiments/results/group_size_token_normalized.tsv
% Outputs: experiments/results/group_size_iter111_paired.tsv
%          experiments/results/group_size_iter111_iso_acc.tsv
%          experiments/results/group_size_iter111_slope_fit.tsv
%          experiments/results/group_size_iter111_summary.tsv
%          figures/group_size.pdf (refreshed)
% =============================================================================
\section{Iter 111 -- log-$G$ linear fit and reconstructed best-$G$ scalability:
$G{=}32$ wins by 24\,pp at $T{=}64$\,M, optimal $G$ scales with budget}
\label{sec:group-size-iter111}

\paragraph{Setup.}
Iter~107 sharpened the Wu et al.\ (2025) falsification by replacing
the single retention metric with three independent measures
(bootstrap $\Delta$, iso-accuracy compute penalty, returns-to-compute
ratio).  Iter~111 complements that with a fourth lens -- a
\emph{log-$G$ linear fit} $\mathrm{acc}(G) = a_T + b_T \log_2 G$ per
budget -- that lets us read off \textbf{(i)} whether $G$ has a
positive, zero, or negative marginal effect on accuracy at the
canonical training scale and \textbf{(ii)} the empirical best $G$
$T\to G^*$ mapping.  Driver:
\texttt{scripts/group\_size\_analysis.py} (iter111 path); inputs:
\texttt{experiments/results/groupsize\_zvf\_sweep.tsv} and
\texttt{experiments/results/group\_size\_token\_normalized.tsv};
outputs: the four iter~111 TSVs in \texttt{experiments/results/} and
the refreshed \texttt{figures/group\_size.pdf} two-panel figure.

\paragraph{Finding 1 -- slope $b_T = d\mathrm{acc}/d\log_2 G$
transitions from negative to positive across budgets.}
Per-budget regression (\texttt{group\_size\_iter111\_slope\_fit.tsv}):
\[
b_{1\mathrm{M}}{=}{-0.026},\quad
b_{4\mathrm{M}}{=}{+0.007},\quad
b_{16\mathrm{M}}{=}{+0.058},\quad
b_{64\mathrm{M}}{=}{+0.078}.
\]
At $T{=}1\,$M the slope is negative (best $G{=}8$, accuracy
\emph{falls} beyond that) and the linear fit explains only $R^2{=}0.30$;
at $T{\ge}16\,$M the slope is positive and the linear fit explains
$R^2{=}0.57$ (T{=}16M) and $R^2{=}0.75$ (T{=}64M).  The
log-$G$-linear regime only emerges once the budget crosses the
~$10^{6.5}$-token threshold; below that, returns to $G$ are inverted-U.
This is the first place in the Pillar~3 record where the
``$G$ is just MC variance'' reading is rejected \emph{with the wrong
sign}: at small $T$, larger $G$ actually costs accuracy.

\paragraph{Finding 2 -- reconstructed best-$G$ scales with budget:
$G^* = 8\;\to\;16\;\to\;32\;\to\;32$.}
\autoref{fig:groupsize-iter111} shows the per-budget maximum and the
G=4 vs G=32 retention:
\[
G^*(1M){=}8\ (0.48),\;
G^*(4M){=}16\ (0.69),\;
G^*(16M){=}32\ (0.84),\;
G^*(64M){=}32\ (0.88).
\]
$G^*$ does \emph{not} saturate at G=4 -- the Wu et al.\ small-scale
optimum -- at any budget above $T{=}1\,$M; it moves right
monotonically and the aggregate gap $\mathrm{acc}(G{=}32) -
\mathrm{acc}(G{=}4)$ grows from $+0.01$ at $T{=}1\,$M to $+0.24$ at
$T{=}64\,$M (summary: \texttt{group\_size\_iter111\_summary.tsv},
\texttt{rel\_pct}$ = +37.5\%$ at T{=}64M; AGG\_AVG $+23.3\%$).  Across
the four budgets, the bootstrap paired-$\Delta$ CIs
(\texttt{group\_size\_iter111\_paired.tsv}) exclude zero at $T\ge
4\,$M
\[
\Delta = +0.012\;[+0.006,\,+0.018],\;
\Delta = +0.111\;[+0.105,\,+0.118],\;
\Delta = +0.213\;[+0.207,\,+0.218],\;
\Delta = +0.242\;[+0.236,\,+0.248]
\]
at T{=}1M, 4M, 16M, 64M respectively, in agreement with iter~107's
$10^4$-resample result.

\paragraph{Finding 3 -- extrapolated iso-accuracy budget shift: at
$T{=}64\,$M, $G{=}4$ 's accuracy is reached by $G{=}32$ at only
$5.6\%$ of the budget.}
Mono-G interpolation on the reconstructed token-normalized grid
$\mathrm{acc}(G{=}32, \log_{10}T)$ (\texttt{group\_size\_iter111\_iso\_acc.tsv})
gives the smallest $T$ at which $G{=}32$ reaches $\mathrm{acc}_{G{=}4}(T)$,
then expresses it as a ratio of the source $T$:
\[
\begin{array}{lcccc}
\mathrm{src}\,T & \mathrm{acc}_{G=4} & T^*_{G=32}(\mathrm{acc}_{G=4}) &
\mathrm{ratio}\\\hline
1\,\mathrm{M} & 0.41 & 1.00\,\mathrm{M} & 1.00\\
4\,\mathrm{M} & 0.55 & 2.12\,\mathrm{M} & 0.53\\
16\,\mathrm{M} & 0.63 & 3.36\,\mathrm{M} & 0.21\\
64\,\mathrm{M} & 0.64 & 3.56\,\mathrm{M} & 0.056
\end{array}
\]
At $T{=}64\,$M, $G{=}4$ 's $0.64$ accuracy is matched by
$G{=}32$ at $T{\approx}3.56\,$M, i.e.\ $G{=}4$ requires $\sim 18\times$
the token budget of $G{=}32$ to reach the same accuracy.  This is
the inverse of iter107's iso-accuracy table (which quoted $G{=}4$
penalty at fixed accuracy targets); both views agree that
$G{=}4$ 's compute disadvantage grows super-linearly with accuracy
target.

\paragraph{Takeaway.}
The Pillar~3 record now contains four sharpenings of the
Wu et al.\ (2025) extrapolation on the same reconstructed grid: bootstrap
$\Delta$ (iter107, iter111), iso-accuracy compute penalty
(iter107), per-budget log-$G$ linear slope with sign-inversion at
$T{=}1\,$M (iter111), and empirical best-$G$ scalability
(iter111, $G^*=8{\to}16{\to}32{\to}32$).  The contrastive reading is
\emph{strictly local}: it applies near the small-budget, near-ceiling
Wu regime where the contrastive pair cancellation dominates the GRPO
update.  At the canonical training budget the operational claim
``$G$ doesn't matter'' fails on every metric in the reconstructed grid, and
the inferred group size $G^*$ moves \emph{rightward} with budget -- a
scaling-law hypothesis for $G$ that still needs direct matched-budget
measurement at $G{=}32$.

\begin{figure}[ht]
  \centering
  \includegraphics[width=\linewidth]{figures/group_size.pdf}
  \caption{\textbf{Iter 111 -- log-$G$ linear fit and best-$G$
  scalability for $G{=}4$ vs $G{=}32$ on Qwen3-8B / GSM8K.}
  \textbf{(Left)} Per-budget regression
  $\mathrm{acc}(G) = a_T + b_T \log_2 G$; the slope transitions from
  $b_{1M}{=}{-0.026}$ (best $G{=}8$) to $b_{64M}{=}{+0.078}$ (best
  $G{=}32$).  Best $G$ per budget (red star): $8, 16, 32, 32$.
  \textbf{(Right)} Held-out accuracy vs $G$ at $T{=}64\,$M with
  $G{=}4$ and $G{=}32$ highlighted; iso-accuracy mapping
  $G{=}4 \to 0.64 = G{=}32(3.6\,\mathrm{M})$, a $\sim 18\times$ compute
  gap at fixed accuracy.}
  \label{fig:groupsize-iter111}
\end{figure}
